\documentclass[a4paper, 12pt]{article}

% =======================
% IMPORTATION DES PACKAGES
% =======================

\usepackage[T1]{fontenc}
\usepackage[utf8]{inputenc}
\usepackage[margin=1.5cm]{geometry}
\usepackage{lmodern}
\usepackage{theorem}
\usepackage{amsfonts, amsmath, amssymb}
\usepackage[french]{babel}

\renewcommand{\familydefault}{\sfdefault}

\title{TOPOLOGIE}
\date{}
\author{}

\theoremstyle{break}
\theoremheaderfont{\scshape}
\theorembodyfont{\upshape}

\newtheorem{defin}{Définition}[section]
\newtheorem{prop}[defin]{Proposition}
\newtheorem{theo}[defin]{Théorème}
\newtheorem{exe}[defin]{Exemple}

\begin{document}

\maketitle

Ceci n'est pas a priori un book, ces sont justes des notes.

\section{Espace topologique}
\begin{defin}
Soit X un espace vectoriel normé et $\tau$ un famille d'ouvert qui contient les partis de X.
On dit que $\tau$ est un topologie si elle vérifie les conditions suivants:
\begin{enumerate}
	\item $\emptyset \in \tau$  et $X \in \tau$
	\item Si $(O_{i})_{i \in \mathbb{N}} \in \tau$, on a $\bigcup_{i \in \mathbb{N}} O_{i} \in \tau$
	\item Si $(O_{i})_{i \in \mathbb{N}} \in \tau$, on a $\bigcap_{i \in \mathbb{N}} O_{i} \in \tau$
\end{enumerate}
\end{defin}
Le couple (X, $\tau$) est un espace topologique. $\tau$ est un ouvert.

\begin{exe}
	\begin{enumerate}
		\item $\tau = P(X)$ est une topologie, elle est appelé topologie discrète.
		\item $\tau = \{0, X\}$ est une topologie appelé topologie grossière.
		\item $\tau = \{ O \subset X, X \backslash O \mbox{ finie}\} \cap \{ \emptyset \}$ est une topologie appelé topologie cofinie.
	\end{enumerate}
\end{exe}
\end{document}
